% ==============================================================================
% Projeto de Sistema - Nome do Aluno
% Capítulo 2 - Plataforma de Desenvolvimento
% ==============================================================================
\chapter{Plataforma de Desenvolvimento}
\label{sec-plataforma}
\vspace{-1cm}


% \vitor{Esta seção deve apresentar a plataforma de implementação a ser adotada para o desenvolvimento do sistema, incluindo: linguagem de programação, mecanismo de persistência de dados e componentes ou \textit{frameworks} a serem usados. Adicionar às tabelas quantas linhas forem necessárias. Remover a marcação em amarelo (comando \textbackslash hl\{\} no \LaTeX). Como exemplo, vide documento de projeto do Marvin. Remover o \hl{fundo amarelo}.}


%=======================================================================================================
%			Tabela de Plataforma de Desenvolvimento e Tecnologias Utilizadas
%=======================================================================================================

Na Tabela~\ref{tabela-plataforma} são listadas as tecnologias utilizadas no desenvolvimento da ferramenta, bem como o propósito de sua utilização.

\begin{footnotesize}
\begin{longtable}{|p{2.5cm}|c|p{5cm}|p{5.5cm}|}
	\caption{Plataforma de Desenvolvimento e Tecnologias Utilizadas.}	
	\label{tabela-plataforma}\\\hline

	\rowcolor{lightgray}
	\textbf{Tecnologia} & \textbf{Versão} & \textbf{Descrição} & \textbf{Propósito} \\\hline 
	\endfirsthead
	\hline
	\rowcolor{lightgray}
	\textbf{Tecnologia} & \textbf{Versão} & \textbf{Descrição} & \textbf{Propósito} \\\hline 
	\endhead
		
	Gleam & 1.15.4 & Linguagem de Programação Funcional & Compilação de programas tanto para a Beam (máquina virtual do Erlang) quanto para JavaScript \\ \hline
	Lustre & 5.6.0 & Framework para construção de templates HTML & Criação do front-end da aplicação \\ \hline
	Wisp & 2.2.2 & Framework para desenvolvimento web & Roteamento HTTP \\ \hline
	Squirrel & 4.6.0 & Framework type-safe de SQL & Realização de querries em banco de dados \\ \hline
	Glow Auth & 1.0.1 & Framework OAuth 2.0 & Autenticação segura de usuários \\ \hline

\end{longtable}
\end{footnotesize}






%=======================================================================================================
%			Tabela de Softwares de Apoio ao Desenvolvimento do Projeto
%=======================================================================================================

Na Tabela~\ref{tabela-software} vemos os softwares que apoiaram o desenvolvimento de documentos e também do código fonte.

\begin{footnotesize}
\begin{longtable}{|p{2.5cm}|c|p{5cm}|p{5.5cm}|}
	\caption{Softwares de Apoio ao Desenvolvimento do Projeto}	
	\label{tabela-software}\\\hline
	
	\rowcolor{lightgray}
	\textbf{Tecnologia} & \textbf{Versão} & \textbf{Descrição} & \textbf{Propósito} \\\hline 
	\endfirsthead
	\hline
	\rowcolor{lightgray}
	\textbf{Tecnologia} & \textbf{Versão} & \textbf{Descrição} & \textbf{Propósito} \\\hline 
	\endhead
		
	TeXstudio & 4.9.3 & Editor LaTeX & Escrita deste documento. \\ \hline
	Visual Studio Code & 1.116.0 & Editor de Código & Escrita do código fonte. \\ \hline
	Mermaid Live Editor & 11.14.0 & Editor de Diagramas & Criação dos diagramas UML. \\ \hline
\end{longtable}
\end{footnotesize}
